% !PRINTED 110
\weberpage{110}
%<<CONT
and in the same way we obtain

\weq{14}{(\varrho^2, x_0) = (-\varepsilon^2)^\lambda \varrho_2 (-\sqrt{-3})^u q_1^{\sigma_1} q_2^{\sigma_2} \ldots p_1^{-\tau_1} p_2^{-\tau_2} \ldots H_2.} % [sic: Weber prints exponent u; the conjugate of (13) requires \mu.]

Because the product of these two expressions must be rational, there results
further, between the roots of unity $\varrho, \varrho_1, \varrho_2$, the relation

\[
\varrho^\lambda \varrho_1 \varrho_2 = 1.
\]

If we set

\weq{15}{x_0 + x_1 + x_2 = A,}

then $A$ too is a rational number, and from (13), (14), (15) it follows that

\weq{16}{\begin{aligned}
x_0 &= 1/3 \left[A + (\varrho, x_0) + (\varrho^2, x_0)\right] \\
x_1 &= 1/3 \left[A + \varrho^2 (\varrho, x_0) + \varrho (\varrho^2, x_0)\right] \\
x_2 &= 1/3 \left[A + \varrho (\varrho, x_0) + \varrho^2 (\varrho^2, x_0)\right]
\end{aligned}}

are cyclotomic numbers. % [sic] printed "Kreistheilungzahlen", elsewhere on the page "Kreistheilungszahlen"
From the formulae (13), (14) we can see the composition of these numbers out of
the periods $\eta_1, \eta_2, \ldots$ But a closer examination of this
subject is no longer required, because already in \S.\,23 we have completely
constructed all cyclotomic numbers that satisfy cubic equations.

\noindent\hrulefill

\begin{center}
\S.\,26.

\gesperrt{Biquadratic Abelian Equations.}
\end{center}

The proof of the corresponding theorem for the biquadratic \gesperrt{Abelian} equations
can be carried out quite similarly. Only here two cases are to be
distinguished, namely equations with non-cyclic group, and
equations with cyclic group.

For the non-cyclic \gesperrt{Abelian} equations with the roots $x_0, x_1, x_2,
x_3$ the group is:

\weq{1}{1, \ \ (0,1)(2,3), \ \ (0,2)(1,3), \ \ (0,3)(1,2),}

and thus the three squares are

\weq{2}{\begin{aligned}
(x_0+x_1-x_2-x_3)^2 &= a \\
(x_0-x_1+x_2-x_3)^2 &= b \\
(x_0-x_1-x_2+x_3)^2 &= c
\end{aligned}}

and the product

\weq{3}{(x_0+x_1-x_2-x_3)(x_0-x_1+x_2-x_3)(x_0-x_1-x_2+x_3) = c,} % [sic: Weber prints RHS c, reusing c from the third square in (2); this rational product requires a distinct symbol.]
%CONT>>
